\section{Version 4.2 Follow-up Studies}
\label{app:v4p2-followups}

This appendix records targeted checks requested after completion of the v4.2
result.  They retain the v4.2 signal model, conversion, physical windows,
statistical conventions, and interpretation boundaries without altering the
accepted v4.2 result.  They are not used to replace the principal observed scan
or its conditional fixed-state bands.

\subsection{Common-binning and profiled 65 MeV extraction}
\label{app:v4p2-common-binning}

The native post-rebin widths in Figure~\ref{fig:results-v4p2-m065-extraction}
are 0.25, 0.25, and 0.625~MeV for 2015, 2016, and 2021, respectively.  The
underlying 2015 and 2016 histogram bins are 0.05~MeV wide, whereas the 2021
source bins are 0.125~MeV wide.  Consequently, an exact common 0.625~MeV
aggregation would require the noninteger factor 12.5 for 2015 and 2016.  The
validated ROOT inputs contain histograms but no event-level tree from which a
new edge-aligned spectrum can be filled.  Splitting, weighting, or rounding
source-bin counts would introduce fractional or reassigned pseudo-counts into
the integer-Poisson likelihood, so that construction is not used.

The exact 0.5 and 0.75~MeV source-compatible choices are equally distant from
0.625~MeV.  The study uses 0.5~MeV because it is the finer choice and the
nearest one that retains the complete support of all three inputs, with integer
aggregation factors 10, 10, and 4; 0.75~MeV would require endpoint trimming in
2015 and 2021.  A second exact 1.25~MeV aggregation with factors 25, 25, and 10
provides a coarser stress test.  Each alternative is a newly optimized
sideband GP refit at 65~MeV using the v4.2 physical resolution, training
exclusion, support, radiative conversion, and profile likelihood.  The count
density entering the signal conversion remains calculated from the uncropped
fine histogram in the physical $m\pm1.64\sigma_m$ window.

\begin{table}[H]
\centering
\small
\caption{Fixed-mass 65~MeV common-binning check.  The quoted uncertainty is the
local signed-estimator Wald curvature scale.  The 1.25~MeV 2015 support ends at
134 rather than 135~MeV because the original support length is not divisible
by 1.25~MeV.  No tolerance for ``unchanged'' significance was declared in
advance, so the shifts are reported directly.}
\label{tab:v4p2-common-binning-m065}
\begin{tabular}{lcccccc}
\toprule
Binning & factors $(15,16,21)$ &
$\widehat{\eps^2}$ & $\sigma_{\eps^2}$ & local $p_0$ & local $Z$ &
$\Delta Z$ \\
 & & \multicolumn{2}{c}{($\times10^{-6}$)} & & & \\
\midrule
Native v4.2 & $(5,5,5)$ & 6.4271 & 1.6101 &
$3.2592\times10^{-5}$ & 3.9932 & 0 \\
Common 0.5~MeV & $(10,10,4)$ & 6.2072 & 1.6110 &
$5.8033\times10^{-5}$ & 3.8543 & $-0.1389$ \\
Common 1.25~MeV & $(25,25,10)$ & 6.4049 & 1.6346 &
$4.4447\times10^{-5}$ & 3.9191 & $-0.0742$ \\
\bottomrule
\end{tabular}
\end{table}

The exact 0.5~MeV result therefore produces a modest fixed-mass shift rather
than numerical identity: the local excess remains near $4\sigma$, while
$Z_{\mathrm{local}}$ decreases by 0.139 and the fitted coupling changes by
$-0.220\times10^{-6}$.  The Wald scale changes by less than
$0.001\times10^{-6}$.  Both independent optimizer-repeat sets select stable,
non-boundary solutions.  The 1.25~MeV stress test gives the same qualitative
conclusion, subject to its recorded one-MeV far-side support trim.

\begin{figure}[H]
\centering
\graphicorplaceholder{0.99\linewidth}{final_limit_projection_figs/v4p2_followups_20260806/figure61_common_0p5MeV.pdf}
\caption{Observed 65~MeV extraction after an exact, count-preserving common-bin
refit.  All three histograms use 0.5~MeV bins.  The upper panels show data, the
newly optimized sideband-trained GP mean and predictive uncertainty, and
display extensions of the standalone and simultaneous shared-$\eps^2$ signal
shapes.  The lower panels show data minus the pre-profile GP mean.  This is a
local binning-robustness check, not a replacement for the accepted native v4.2
scan.  Its lower-panel error bars use the prior-predictive scale
$\sqrt{\mu_{\mathrm{GP}}+C_{ii}}$ and are not independent post-fit bin
uncertainties.}
\label{fig:appendix-v4p2-common-binning-m065}
\end{figure}

\begin{figure}[H]
\centering
\graphicorplaceholder{0.99\linewidth}{final_limit_projection_figs/v4p2_followups_20260806/figure61_common_0p5MeV_profiled.pdf}
\caption{Profiled-background view of the exact common-0.5~MeV extraction,
restricted to the actual $\pm2.25\sigma_m$ likelihood window.  Blue denotes
the background-only profiled expectation, orange the standalone
signal-plus-background profile, and red the simultaneous shared-$\eps^2$
profile.  The lower panels show count residuals relative to each expectation.
The curves are not extended into the sidebands because the v4.2 nuisance
likelihood profiles the GP background only inside the extraction window.
These correlated fit diagnostics are not standardized per-bin significances.}
\label{fig:appendix-v4p2-common-binning-profiled}
\end{figure}

The negative 2016 endpoint in the original lower-right panel of
Figure~\ref{fig:results-v4p2-m065-residuals} came from plotting the signed,
unconstrained estimate
$(1.51921\pm2.55749)\times10^{-6}$ with a symmetric Wald error.  Its lower
endpoint is $-1.03828\times10^{-6}$ even though the event-to-$\eps^2$
conversion factor is positive.  It is therefore not a sign or normalization
error, but it is easy to misread as a physical interval.  The corrected
display uses the physical domain $\eps^2\geq0$ and the nominal asymptotic
profile-likelihood set defined by $\Delta(-2\ln L)=1$.  For native 2016 that
set is $[0,4.08001]\times10^{-6}$; it is not asserted to be
coverage-calibrated.

\begin{figure}[H]
\centering
\graphicorplaceholder{0.92\linewidth}{final_limit_projection_figs/v4p2_followups_20260806/figure62_coefficients_physical68.pdf}
\caption{Physical-domain 65~MeV signal-strength estimates for the native v4.2
extraction and the exact common-0.5~MeV refit.  Horizontal bars are the nominal
asymptotic $\eps^2\geq0$ profile sets satisfying
$\Delta(-2\ln L)\leq1$.  The 2016 lower endpoints are zero.  These sets replace
the ambiguous symmetric signed-Wald display; they are not
coverage-calibrated confidence intervals.}
\label{fig:appendix-v4p2-physical-profile-intervals}
\end{figure}

\subsection{Conditional 65 MeV central-window pseudo-observation study}
\label{app:v4p2-pseudo65}

This study replaces the 65~MeV central region of the 2021 10\% spectrum with
two smooth-background pseudo-observations and then treats each replacement as
an observed dataset.  The requested continuous interval is
$65~\mathrm{MeV}\pm2.5\sigma_m=[59.6958,70.3042]$~MeV, using
$\sigma_m(65~\mathrm{MeV})=2.121696875$~MeV.  On the production 0.625~MeV
grid, the $\pm2.25\sigma_m$ and $\pm2.5\sigma_m$ center selections happen to
select the same sixteen analysis bins, whose complete edges are
$[60,70)$~MeV.  The ROOT construction therefore replaces the corresponding
80 native 0.125~MeV bins.  All bins outside this interval are bitwise identical
to the original observed spectrum.

No nonzero signal strength was supplied, so both lanes use the fail-closed
background-only interpretation $A_{\mathrm{inj}}=0$.  The first lane draws
independent binwise Poisson counts from the exact accepted v4.2 fixed-GP mean
at 65~MeV.  The second draws from a sideband-only
\texttt{fGenGammaThresh} interpolation fitted with a binned Poisson
likelihood over 50--85~MeV with $[60,70)$~MeV excluded.  Five deterministic
starts were compared.  The selected functional fit has Poisson
deviance/ndf 1.069, a deviance $p$-value of 0.243, and no selected parameter
at its declared bound.  It is a smooth interpolation stress truth, not a
physical event generator.  Independent PCG64 child streams of master seed
20260806 generate the two draws.  The source copy, both pseudo-observed
histograms, and both central expectations are written to one ROOT file with
their construction and array hashes.

\begin{figure}[H]
\centering
\graphicorplaceholder{0.99\linewidth}{final_limit_projection_figs/v4p2_followups_20260806/pseudo65_central_window_zoom.pdf}
\caption{Construction of the two fixed-seed background-only replacements in
the 65~MeV central window.  Gray shows the original 2021 10\% data, blue or
orange the integer-valued conditional pseudo-observation, and the dashed black
curve the corresponding replacement expectation.  The shaded interval has
complete-bin edges $[60,70)$~MeV and the dotted line marks 65~MeV.  The two
panels are single Poisson draws, not ensemble means; only the outside-window
data and the bin geometry, not the replaced total count, are preserved.}
\label{fig:appendix-v4p2-pseudo65-construction}
\end{figure}

Each pseudo-observation is passed through the full 2021-only v4.2
observed/asymptotic scan from 50 to 250~MeV: 0.625~MeV analysis bins,
$\pm2.25\sigma_m$ extraction and training geometry, $k_{\max}=15$, profiled
correlated-background extraction, 90\% \CLs{}, and a local bounded
asymptotic $p_0$.  Thus the wider requested interval controls the complete
central bins replaced in the ROOT file but does not redefine the inference
card.  No expected limit bands or toys are produced.

Two unchanged-card full scans were compared for each lane.  Where they selected
different optimizer branches, isolated mass points were rerun with the
unchanged card and the reproducible maximum-log-marginal-likelihood state was
selected; no point was interpolated.  The final ledgers contain all 201 masses
per lane, have no pending or nonfinite results, and select no length-scale
boundary states.  Multiple finite branches were recorded at 12 GP-mean masses
and 11 functional-form masses rather than hidden.

\begin{table}[H]
\centering
\small
\caption{Observed/asymptotic 2021-only results at 65~MeV.  The original v4.2
row is included only as context.  Negative replacement estimates are allowed
in the signed local fit; their bounded one-sided discovery statistic gives
$p_0=0.5$, which is not a posterior probability that the background model is
true.}
\label{tab:v4p2-pseudo65-m065}
\begin{tabular}{lrrrrr}
\toprule
Sample & $\widehat A$ & $\sigma_A$ &
$\eps^2_{90}$ & local $p_0$ & local $Z$ \\
 & \multicolumn{2}{c}{events} & $(\times10^{-6})$ & & \\
\midrule
Original 2021 10\% & 28038.9 & 6609.5 & 11.7184 &
$1.0570\times10^{-5}$ & 4.2525 \\
GP-mean replacement & $-5125.8$ & 6602.8 & 2.7736 &
0.5 & 0 \\
Functional-form replacement & $-2195.6$ & 6603.0 & 3.1242 &
0.5 & 0 \\
\bottomrule
\end{tabular}
\end{table}

\begin{figure}[H]
\centering
\graphicorplaceholder{0.99\linewidth}{final_limit_projection_figs/v4p2_followups_20260806/pseudo65_observed_limit_p0_aligned.pdf}
\caption{Conditional 2021 10\% replacement scans.  The columns use the
GP-mean and functional-form draws, respectively.  Within each column the
pseudo-observed spectrum, observed 90\% \CLs{} upper limit on $\eps^2$, and
local asymptotic $p_0$ are aligned in mass.  The original v4.2 observed curves
are dashed context only; no expected bands are constructed.  At 65~MeV the
replacement lanes give $\eps^2_{90}=2.77357\times10^{-6}$ and
$3.12419\times10^{-6}$, respectively, with $p_0=0.5$ in both cases.  The
tighter single-draw limits are fluctuation dependent and are not expected
sensitivities.}
\label{fig:appendix-v4p2-pseudo65-scans}
\end{figure}

At 65~MeV, the localized positive estimate in the original 2021 10\% spectrum
is absent in both fixed replacement draws.  Across the full conditional scans,
however, the smallest local values occur at the lower 50~MeV scan boundary:
$p_0=1.51\times10^{-4}$ ($Z=3.61$) for the GP-mean lane and
$p_0=9.81\times10^{-7}$ ($Z=4.76$) for the functional-form lane.  Replacing
$[60,70)$~MeV can change the full-support sideband-trained GP used at other
hypotheses, while the outside-window data remain observed.  The boundary
minima are therefore reported as conditional fixed-mass diagnostics and are
not scan calibrated or globalized.  The disappearance of the 65~MeV estimate
is likewise a conditional counterfactual observation, not a probability
statement about the source of the feature.  Because both lanes retain the
same observed spectrum outside the replacement window and contain only one
draw each, they are not independent global-null pseudoexperiments and do not
establish expected sensitivity, coverage, or a global discovery probability.

\subsubsection{Post-Figure~\ref{fig:appendix-v4p2-pseudo65-scans}
functional-form shoulder check}

The functional-form lane in
Figure~\ref{fig:appendix-v4p2-pseudo65-scans} has a conspicuous local shoulder
below 65~MeV.  Its fixed-seed draw gives
$p_0=1.12297\times10^{-4}$ ($Z=3.6896$) at 62~MeV, with
$\widehat A=23177.5$ and $\sigma_A=6288.4$ events.  To separate the smooth
truth shape from this particular Poisson fluctuation, the central
$[60,70)$~MeV bins were replaced by the stored fractional functional mean
itself and the unchanged v4.2 analysis was repeated over 55--75~MeV.
The corresponding fixed-GP-mean construction was analyzed in parallel.  These
are deterministic central-mean diagnostics with the observed spectrum retained
outside the replacement interval; they are not complete Asimov datasets.

The functional mean alone produces a broad positive response: it peaks at
$Z=2.268$ at 61~MeV and gives $Z=1.909$ at 62~MeV.  The fixed-GP mean gives
$Z=0.951$ and 0.431 at the same masses.  At 62~MeV the end-to-end signed
functional-mean response exceeds the fixed-GP-mean response by 9369.2 fitted
events, about $1.47\sigma_A$.  The particular functional Poisson draw adds a
further 11062.3 fitted events relative to its deterministic mean response,
about $1.76\sigma_A$, and moves the local maximum from 61 to 62~MeV.  Thus the
Figure~\ref{fig:appendix-v4p2-pseudo65-scans} shoulder contains both a coherent
functional-truth/GP-analysis shape-mismatch component and an additional
realized fluctuation; it should not be read as a feature of the observed 2021
data.

\begin{figure}[H]
\centering
\graphicorplaceholder{0.92\linewidth}{final_limit_projection_figs/v4p2_followups_20260806/functional_mean_shape_bias_Ahat_p0.pdf}
\caption{Conditional central-mean shape diagnostic over 55--75~MeV.  Solid
curves replace $[60,70)$~MeV by the stored fractional functional or fixed-GP
expectation; dashed curves are the corresponding fixed-seed Poisson draws from
Figure~\ref{fig:appendix-v4p2-pseudo65-construction}.  The upper panel shows
the signed fitted signal yield and the lower panel the local asymptotic
$p_0$.  The smooth functional mean already produces a 61--63~MeV shoulder,
while the displayed Poisson draw amplifies it.  All 42 selected optimizer
states were reproduced by unchanged-card repeats, with no selected
length-scale boundary and no interpolation.  This is a conditional
truth-model/analysis-model diagnostic, not a calibrated bias, expected
sensitivity, or global-significance result.}
\label{fig:appendix-v4p2-functional-mean-shape}
\end{figure}

The functional interpolation nevertheless passes its declared sideband
goodness-of-fit checks: Poisson deviance/ndf is 1.069, the deviance
$p$-value is 0.243, five starts reach the near-optimal solution, and no
selected parameter lies on its bound.  Those tests assess the interpolation
on its fitted sidebands; they do not guarantee zero signed signal response
when a moving-mask GP analyzes that smooth truth.

\subsubsection{GP-mean replacement-window ensembles}

At 65~MeV the requested replacement half-widths have the following exact
production-grid interpretation:

\begin{table}[H]
\centering
\small
\caption{Continuous replacement intervals and the complete 0.625~MeV
analysis bins selected by their centers.  The 2.25- and 2.5-$\sigma_m$
requests are exactly the same binned operation and therefore cannot define
independent window-scope results.}
\label{tab:v4p2-pseudo65-window-geometry}
\begin{tabular}{lccc}
\toprule
Half-width & Continuous interval [MeV] & Complete-bin edges [MeV] &
analysis bins \\
\midrule
$2.25\sigma_m$ & $[60.2262,69.7738]$ & $[60.00,70.00)$ & 16 \\
$2.5\sigma_m$  & $[59.6958,70.3042]$ & $[60.00,70.00)$ & 16 \\
$3.0\sigma_m$  & $[58.6349,71.3651]$ & $[58.75,71.25)$ & 20 \\
\bottomrule
\end{tabular}
\end{table}

Ten independent background-only Poisson replacements were generated from the
exact accepted fixed-GP mean for the shared 2.25/2.5-$\sigma_m$ geometry and
ten for the 3-$\sigma_m$ geometry.  Matching draw indices use common random
numbers: their counts are identical in the sixteen shared bins, while the
wider construction adds four replaced edge bins.  The inference card itself
remains the frozen v4.2 2.25-$\sigma_m$ card, so this comparison isolates
replacement scope rather than simultaneously changing the likelihood
geometry.  Each histogram retains the observed spectrum outside its declared
replacement window.

\begin{figure}[H]
\centering
\graphicorplaceholder{0.99\linewidth}{final_limit_projection_figs/v4p2_followups_20260806/gp_window_ensemble_central_spectra.pdf}
\caption{The ten paired fixed-GP-mean pseudo-observations near 65~MeV.
The left column uses the shared 2.25/2.5-$\sigma_m$ sixteen-bin replacement
and the right column the 3-$\sigma_m$ twenty-bin replacement.  Upper panels
show the individual integer draws, their arithmetic mean, the fixed
generating mean, and the retained original data outside the replaced bins.
Lower panels show each draw's Poisson residual relative to the generating
mean.  Matching draw indices contain identical counts in the shared
$[60,70)$~MeV interval.}
\label{fig:appendix-v4p2-gp-window-spectra}
\end{figure}

\begin{table}[H]
\centering
\small
\caption{Pointwise 65~MeV summaries across ten conditional fixed-GP-mean
draws.  The 16--84\% intervals are empirical descriptive spreads, not
expected-limit bands.}
\label{tab:v4p2-pseudo65-gp-ensemble-m065}
\begin{tabular}{lccccc}
\toprule
Replacement geometry & median $\eps^2_{90}$ & mean $\eps^2_{90}$ &
empirical 16--84\% & median $p_0$ & median $Z$ \\
 & \multicolumn{3}{c}{$(\times10^{-6})$} & & \\
\midrule
$2.25\sigma_m=2.5\sigma_m$ & 3.69710 & 3.57686 &
$[2.70526,4.14840]$ & 0.437686 & 0.1569 \\
$3.0\sigma_m$ & 4.27450 & 4.19152 &
$[3.06747,5.12518]$ & 0.284313 & 0.5702 \\
\bottomrule
\end{tabular}
\end{table}

\begin{figure}[H]
\centering
\graphicorplaceholder{0.99\linewidth}{final_limit_projection_figs/v4p2_followups_20260806/gp_window_ensemble_observed_limit_p0.pdf}
\caption{Dedicated fixed-GP-mean replacement ensembles.  The columns compare
the shared 2.25/2.5-$\sigma_m$ and 3-$\sigma_m$ replacement scopes.  The rows
show the mass spectrum, observed 90\% \CLs{} upper limit, and local
asymptotic $p_0$, aligned in mass.  Thin curves are the ten individual
conditional draws; thick colored and dashed black curves are their pointwise
median and arithmetic mean, and the translucent region is the empirical
16--84\% spread.  The original observed v4.2 curves are gray dashed context.
The empirical spread is not an expected-limit band, and the ensemble is not
a scan-wide null calibration.}
\label{fig:appendix-v4p2-gp-window-ensemble}
\end{figure}

At 65~MeV the paired median ratio
$\eps^2_{90}(3\sigma_m)/\eps^2_{90}(2.25/2.5\sigma_m)$ is 1.17957, with
empirical 16--84\% spread $[1.00971,1.27080]$.  The normalization density is
identical for each paired draw because its physical
$m\pm1.64\sigma_m$ window lies wholly inside the shared replaced bins; the
difference comes from the four outer bins and the profiled GP/extraction
response.  None of the twenty conditional draws reaches local $3\sigma$ in
55--75~MeV.  The smallest values are $p_0=0.00692$ ($Z=2.462$) for the
narrow geometry and $p_0=0.01118$ ($Z=2.284$) for the wide geometry, both at
60~MeV in draw 04.

All 420 rows in the predeclared 55--75~MeV review interval have a reproduced
maximum-finite-log-marginal-likelihood selected state after unchanged-card
repeats.  Seventeen rows retain documented multiple finite branches; none has
an unresolved selected state, selected kernel bound, or interpolated result.
Outside that interval the displayed full-range curves have one scan attempt
with twelve within-fit restarts, so full-grid repeat stability is not claimed.
These ten-draw summaries are a conditional recurrence and mechanism screen,
not expected sensitivity, coverage, or a probability for a future full-data
observation.
